% !TEX program = pdflatex
\documentclass[aspectratio=169]{beamer}

\usetheme{Madrid}
\usecolortheme{default}

\usepackage{amsmath,amssymb,mathtools}
\usepackage{physics}

\title{Measurement Strategies in Quantum Computing}
\subtitle{Why Measuring All Qubits is Optimal in Modern Algorithms}
\author{Morten Hjorth-Jensen}
\date{Spring 2026}

\begin{document}

\begin{frame}
\titlepage
\end{frame}

%------------------------------------------------

\begin{frame}{Motivation}
\begin{block}{Central question}
Why is it preferable to measure all qubits rather than mapping observables onto a single qubit?
\end{block}

\begin{itemize}
\item Traditional textbooks emphasize single-qubit readout.
\item Modern quantum algorithms (VQE, QML) rely on full-register measurement.
\end{itemize}
\end{frame}

%------------------------------------------------

\begin{frame}{Quantum measurement structure}
An $n$-qubit state is
\[
|\psi\rangle = \sum_{x \in \{0,1\}^n} c_x |x\rangle.
\]

Measurement in the computational basis yields:
\[
p(x) = |c_x|^2.
\]

\medskip

\begin{block}{Key observation}
The outcome is a \textbf{joint probability distribution} over all qubits.
\end{block}
\end{frame}

%------------------------------------------------

\begin{frame}{Theorem: Information loss in partial measurement}
\begin{theorem}
Let $\rho$ be an $n$-qubit state. Measuring only a subset of qubits yields marginal distributions which do not uniquely determine $\rho$ or its multi-qubit correlations.
\end{theorem}

\begin{proof}
Let $\rho$ be entangled:
\[
\rho \neq \rho_1 \otimes \rho_2.
\]

Then in general:
\[
\langle Z_i Z_j \rangle \neq \langle Z_i \rangle \langle Z_j \rangle.
\]

Measuring only qubit $i$ gives access to $\langle Z_i \rangle$, but not to joint correlators.

Hence the reduced statistics are insufficient to reconstruct expectation values of multi-qubit observables.
\end{proof}
\end{frame}

%------------------------------------------------

\begin{frame}{Corollary: Necessity of full measurement}
\begin{block}{Corollary}
To estimate expectation values of general observables,
\[
H = \sum_k c_k P_k,
\]
with $P_k$ multi-qubit operators, one must measure all qubits involved.
\end{block}

\[
\langle P_k \rangle = \langle Z_{i_1} Z_{i_2} \cdots X_{j_1} \rangle
\]

\medskip

\begin{itemize}
\item Requires joint measurement outcomes
\item Cannot be reconstructed from single-qubit marginals
\end{itemize}
\end{frame}

%------------------------------------------------

\begin{frame}{Theorem: Circuit-depth optimality}
\begin{theorem}
Let $O$ be a multi-qubit observable. Any protocol that maps $O$ onto a single qubit prior to measurement requires a unitary transformation whose circuit depth is at least proportional to the support size of $O$.
\end{theorem}

\begin{proof}[Sketch]
Mapping $O$ to a single qubit requires a unitary $U$ such that
\[
U^\dagger Z_1 U = O.
\]

If $O$ acts on $k$ qubits, then $U$ must entangle these qubits with qubit $1$.

Thus:
\begin{itemize}
\item Requires at least $\mathcal{O}(k)$ entangling gates
\item Increases circuit depth and noise
\end{itemize}

Direct measurement avoids this overhead.
\end{proof}
\end{frame}

%------------------------------------------------

\begin{frame}{Hardware argument}
\begin{block}{Asymmetry in cost}
\begin{itemize}
\item Quantum gates: noisy and expensive
\item Measurements: comparatively cheap and parallelizable
\end{itemize}
\end{block}

\medskip

Thus:
\[
\text{Adding gates (basis rotations, mappings)} \gg \text{measuring more qubits}
\]
\end{frame}

%------------------------------------------------

\begin{frame}{Deferred measurement principle}
\begin{block}{Principle}
Measurements can be deferred until the end of a quantum circuit without changing outcome probabilities.
\end{block}

\[
\Rightarrow \text{One can measure all qubits at the end}
\]

\medskip

This supports the modern paradigm of full-register measurement.
\end{frame}

%------------------------------------------------

\begin{frame}{Modern paradigm (NISQ era)}
\begin{block}{Variational algorithms}
In hybrid algorithms:
\[
\text{prepare state} \rightarrow \text{measure all qubits} \rightarrow \text{classical post-processing}
\]
\end{block}

\medskip

\begin{itemize}
\item Introduced with VQE (2014)
\item Core workflow: sampling full bitstrings
\end{itemize}

\end{frame}


\begin{frame}{Key question}
\begin{block}{Problem}
Can we reconstruct expectation values of a Hamiltonian using only single-qubit measurements?
\end{block}

\medskip

We show:
\begin{itemize}
\item This fails in general
\item It fails even for simple two-qubit systems
\end{itemize}
\end{frame}

%------------------------------------------------

\begin{frame}{Theorem: Failure of marginal reconstruction}
\begin{theorem}
There exist quantum states $\rho$ and observables $H$ such that the expectation value $\langle H \rangle$ cannot be reconstructed from all single-qubit measurement statistics.
\end{theorem}

\begin{proof}
Consider a two-qubit system and the observable
\[
H = Z_1 Z_2.
\]

Take two states:

\[
|\Phi^+\rangle = \frac{1}{\sqrt{2}}(|00\rangle + |11\rangle),
\quad
\rho_{\text{mix}} = \frac{1}{2}(|00\rangle\langle00| + |11\rangle\langle11|).
\]

Both have identical single-qubit marginals:
\[
\langle Z_1 \rangle = \langle Z_2 \rangle = 0.
\]

However:
\[
\langle Z_1 Z_2 \rangle_{\Phi^+} = 1,
\quad
\langle Z_1 Z_2 \rangle_{\text{mix}} = 1.
\]

Now compare with:
\[
\rho' = \frac{1}{2}(|01\rangle\langle01| + |10\rangle\langle10|),
\]

which also has identical marginals but:
\[
\langle Z_1 Z_2 \rangle_{\rho'} = -1.
\]

Thus, marginals do not determine $\langle H \rangle$.
\end{proof}
\end{frame}

%------------------------------------------------

\begin{frame}{Explicit VQE counterexample}
Consider Hamiltonian:
\[
H = Z_1 Z_2.
\]

\medskip

Ansatz:
\[
|\psi(\theta)\rangle = \cos\theta |00\rangle + \sin\theta |11\rangle.
\]

\medskip

Exact energy:
\[
E(\theta) = \langle Z_1 Z_2 \rangle = 1.
\]

\medskip

Single-qubit estimate:
\[
\langle Z_1 \rangle = \cos^2\theta - \sin^2\theta,
\quad
\langle Z_2 \rangle = \cos^2\theta - \sin^2\theta.
\]

Naive reconstruction:
\[
\langle Z_1 Z_2 \rangle \approx \langle Z_1 \rangle \langle Z_2 \rangle
= (\cos^2\theta - \sin^2\theta)^2.
\]

\begin{block}{Failure}
This is \textbf{incorrect} except at special points.
\end{block}
\end{frame}

%------------------------------------------------

\begin{frame}{Error in VQE energy landscape}
True energy:
\[
E_{\text{true}}(\theta) = 1.
\]

Naive single-qubit estimate:
\[
E_{\text{naive}}(\theta) = (\cos 2\theta)^2.
\]

\medskip

\begin{itemize}
\item True landscape: flat
\item Naive landscape: oscillatory
\end{itemize}

\medskip

\begin{block}{Consequence}
Optimization will converge to \textbf{wrong minima}.
\end{block}
\end{frame}

%------------------------------------------------

\begin{frame}{Theorem: VQE failure under marginal approximation}
\begin{theorem}
Replacing multi-qubit correlators with products of single-qubit expectations leads to incorrect energy landscapes and can produce false minima in VQE.
\end{theorem}

\begin{proof}
Let $H = Z_1 Z_2$.
True:
\[
E(\theta) = \langle Z_1 Z_2 \rangle.
\]
Approximation:
\[
E'(\theta) = \langle Z_1 \rangle \langle Z_2 \rangle.
\]

Since correlations are neglected:
\[
E'(\theta) \neq E(\theta),
\]
and the optimization problem changes. Thus the minimizer of $E'$ is not the minimizer of $E$.
\end{proof}
\end{frame}



%------------------------------------------------

\begin{frame}{Why textbooks differ}
\begin{itemize}
\item Early algorithms (Deutsch–Jozsa, phase estimation)
\begin{itemize}
\item Encode answer in a single qubit/register
\end{itemize}

\item Analytical simplifications:
\begin{itemize}
\item Easier to track one observable
\end{itemize}

\item Phase kickback techniques:
\begin{itemize}
\item Transfer information to one qubit
\end{itemize}
\end{itemize}

\medskip

\begin{block}{Conclusion}
These are \textbf{pedagogical or theoretical constructions}, not optimal experimental strategies.
\end{block}
\end{frame}

%------------------------------------------------

\begin{frame}{Final theorem: Optimal measurement strategy}
\begin{theorem}
For variational and sampling-based quantum algorithms, measuring all qubits in the computational basis minimizes circuit depth while preserving full information about the quantum state.
\end{theorem}

\begin{proof}[Argument]
\begin{itemize}
\item Full measurement yields complete probability distribution $p(x)$
\item Multi-qubit observables are computed from bitstrings
\item Avoids additional unitary transformations
\item Reduces noise accumulation
\end{itemize}

Thus, full-register measurement is optimal under realistic hardware constraints.
\end{proof}
\end{frame}


\begin{frame}{Final conclusion}
\begin{itemize}
\item Single-qubit measurements give only marginals
\item Correlations are essential for quantum advantage
\item VQE fails if correlations are ignored
\item Full-register measurement is fundamentally necessary
\end{itemize}

\medskip

\begin{block}{Key insight}
Quantum algorithms operate on \textbf{correlations}, not individual qubits.
\end{block}
\end{frame}



%------------------------------------------------

\begin{frame}{Summary}
\begin{itemize}
\item Quantum information is inherently \textbf{non-local}
\item Single-qubit measurement loses correlations
\item Mapping observables increases circuit depth
\item Modern hardware supports parallel measurement
\item NISQ algorithms rely on full bitstring sampling
\end{itemize}

\medskip

\begin{block}{Key insight}
Quantum computing is fundamentally about sampling high-dimensional distributions, not extracting single bits.
\end{block}
\end{frame}

%------------------------------------------------

\end{document}



%------------------------------------------------

\begin{frame}{Full VQE optimization experiment}
\begin{block}{Goal}
Compare optimization using:
\begin{itemize}
\item True energy: $\langle Z_1 Z_2 \rangle$
\item Marginal approximation: $\langle Z_1 \rangle \langle Z_2 \rangle$
\end{itemize}
\end{block}

\medskip

\textbf{Hamiltonian:}
\[
H = Z_1 Z_2
\]

\textbf{Ansatz:}
\[
|\psi(\theta)\rangle = \cos\theta |00\rangle + \sin\theta |11\rangle
\]
\end{frame}

%------------------------------------------------


\begin{frame}[fragile]{Python: Full VQE loop}
\begin{lstlisting}
import numpy as np

# Pauli matrices
Z = np.array([[1,0],[0,-1]])
Z1Z2 = np.kron(Z, Z)

def state(theta):
    return np.array([np.cos(theta), 0, 0, np.sin(theta)])

def expectation(psi, H):
    return np.real(psi.conj().T @ H @ psi)

def single_qubit_expectations(psi):
    rho = np.outer(psi, psi.conj())
    Z1 = np.kron(Z, np.eye(2))
    Z2 = np.kron(np.eye(2), Z)
    z1 = np.trace(rho @ Z1)
    z2 = np.trace(rho @ Z2)
    return np.real(z1), np.real(z2)
\end{lstlisting}
\end{frame}

%------------------------------------------------

\begin{frame}[fragile]{Gradient descent optimizer}
\begin{lstlisting}
def optimize(cost_fn, theta0, lr=0.1, steps=50):
    theta = theta0
    trajectory = []

    for _ in range(steps):
        # finite difference gradient
        eps = 1e-6
        grad = (cost_fn(theta + eps) - cost_fn(theta - eps)) / (2*eps)

        theta = theta - lr * grad
        trajectory.append((theta, cost_fn(theta)))

    return theta, trajectory
\end{lstlisting}
\end{frame}

%------------------------------------------------

\begin{frame}[fragile]{Define cost functions}
\begin{lstlisting}
# True VQE energy
def true_energy(theta):
    psi = state(theta)
    return expectation(psi, Z1Z2)

# WRONG: marginal approximation
def approx_energy(theta):
    psi = state(theta)
    z1, z2 = single_qubit_expectations(psi)
    return z1 * z2
\end{lstlisting}
\end{frame}

%------------------------------------------------

\begin{frame}[fragile]{Run optimization}
\begin{lstlisting}
theta0 = 0.3

theta_true, traj_true = optimize(true_energy, theta0)
theta_approx, traj_approx = optimize(approx_energy, theta0)

print("True optimum theta:", theta_true)
print("Approx optimum theta:", theta_approx)

print("True energy:", true_energy(theta_true))
print("Approx energy:", approx_energy(theta_approx))
\end{lstlisting}
\end{frame}

%------------------------------------------------

\begin{frame}{Expected outcome}
\begin{itemize}
\item True energy:
\[
E(\theta) = 1 \quad \Rightarrow \text{flat landscape}
\]

\item Approximate energy:
\[
E_{\text{approx}}(\theta) = (\cos 2\theta)^2
\]

\item Minimization gives:
\[
\theta = \frac{\pi}{4}
\]
\end{itemize}

\medskip

\begin{block}{Key result}
Optimizer converges to a \textbf{spurious minimum}.
\end{block}
\end{frame}

%------------------------------------------------

\begin{frame}{Theorem: Wrong convergence in VQE}
\begin{theorem}
If multi-qubit correlations are approximated by products of marginals, the resulting VQE optimization converges to incorrect parameters and energies.
\end{theorem}

\begin{proof}
The approximate energy defines a different cost function:
\[
E'(\theta) \neq E(\theta)
\]

Thus:
\[
\arg\min_\theta E'(\theta) \neq \arg\min_\theta E(\theta)
\]

Since optimization follows $E'$, it converges to the wrong solution.
\end{proof}
\end{frame}

%------------------------------------------------

\begin{frame}{Interpretation}
\begin{itemize}
\item VQE is a \textbf{variational optimization loop}  [oai_citation:0‡grove-docs.readthedocs.io](https://grove-docs.readthedocs.io/en/latest/vqe.html?utm_source=chatgpt.com)
\item The optimizer only sees the measured energy
\item If measurement is wrong $\Rightarrow$ optimization is wrong
\end{itemize}

\medskip

\begin{block}{Deep insight}
Measurement is not a post-processing step — it defines the \textbf{objective function}.
\end{block}
\end{frame}

%------------------------------------------------

\begin{frame}{Final conclusion}
\begin{itemize}
\item Single-qubit measurement destroys correlations
\item Correlations determine energy in VQE
\item Wrong measurement $\Rightarrow$ wrong landscape
\item Wrong landscape $\Rightarrow$ wrong convergence
\end{itemize}

\medskip

\begin{block}{Final statement}
\[
\textbf{Full-register measurement is necessary for correct VQE optimization.}
\]
\end{block}
\end{frame}



%------------------------------------------------

\begin{frame}{Final result: optimizer failure}
\begin{itemize}
\item Correct VQE:
\begin{itemize}
\item Converges to ground state
\item Uses full correlations
\end{itemize}

\item Marginal-based VQE:
\begin{itemize}
\item Converges to wrong state
\item Ignores entanglement
\end{itemize}
\end{itemize}

\medskip

\begin{block}{Fundamental conclusion}
\[
\textbf{Incorrect measurement $\Rightarrow$ incorrect gradient $\Rightarrow$ incorrect physics}
\]
\end{block}
\end{frame}



%------------------------------------------------

\begin{frame}{Final result: optimizer failure}
\begin{itemize}
\item Correct VQE:
\begin{itemize}
\item Converges to ground state
\item Uses full correlations
\end{itemize}

\item Marginal-based VQE:
\begin{itemize}
\item Converges to wrong state
\item Ignores entanglement
\end{itemize}
\end{itemize}

\medskip

\begin{block}{Fundamental conclusion}
\[
\textbf{Incorrect measurement $\Rightarrow$ incorrect gradient $\Rightarrow$ incorrect physics}
\]
\end{block}
\end{frame}
